% Paper C -- working draft, pending revision.
% All claims cite machine verifications in manuflog/contextuality-obstructions.
\documentclass[11pt]{article}
\usepackage{amsmath,amssymb,amsthm,geometry,booktabs,hyperref}
\geometry{margin=1.1in}
\newtheorem{theorem}{Theorem}
\newtheorem{proposition}{Proposition}
\newtheorem{observation}{Observation}
\newtheorem{openproblem}{Open Problem}
\newtheorem{conjecture}{Conjecture}
\newtheorem{lemma}{Lemma}
\newtheorem{theorem2}[theorem]{Theorem}
\newcommand{\U}{\mathrm{U}}
\newcommand{\PU}{\mathrm{PU}}
\newcommand{\avn}{\mathrm{AvN}}
\title{What Contextual Holonomy Detects:\\ the All-versus-Nothing Sector, Two Gap Certificates,\\ and the Section Dependence of Wigner--Friend Holonomy}
\author{Manuel Flores Gordillo\\ \small Independent Researcher \quad ORCID 0009-0006-8246-0343\\ \small \texttt{manuel.flores@columbia.edu}}
\date{July 2026 --- working draft v0}
\begin{document}
\maketitle

\begin{abstract}
The slogan ``contextual holonomy encodes quantum contextuality'' is attractive and, as
usually stated, false. We give it a sharp, machine-verified form. On the context stack
$[X/\PU(n)]$ of maximal abelian subalgebras the canonical degree-2 class has order exactly
$n$; its cocycle holonomy along closed context loops detects \emph{exactly} the
all-versus-nothing (AvN) sector of contextuality for Weyl context families
(Theorem~\ref{thm:equiv}), a two-directional equivalence stress-tested on $2{,}100$ random
families at $d=2,4$ with zero mismatches. Both containments are strict: at $d=3$ the class
is nontrivial of order $3$ while the achievable state-independent AvN value is $0$
(Proposition~\ref{prop:gap1}, exhaustive certificate), and KCBS exhibits state-dependent
contextuality where the class shadow vanishes (Proposition~\ref{prop:gap2}). Finally, the
holonomy of the Wigner--friend loop generated by $H$ and $S$ splits by level: its
projective (Pancharatnam) holonomy is $+1$, while its determinant-section holonomy is
exactly $-i$, a $\tau$-level invariant of the 2-adic tower
(Observation~\ref{obs:wf}). Interpretational claims must therefore specify loop and level.
We state the residual equivalence problem precisely and give hardware-testable predictions.
\end{abstract}

\section{Setting and the repaired slogan}
Let $X$ be the space of maximal abelian $*$-subalgebras (MASAs) of $M_n(\mathbb{C})$,
$n\geq 3$. The unitary group acts by conjugation; on $\U(n)\ltimes X$ the
$\U(1)$-lifting class is vacuous, and after Morita reduction the canonical class on
$\PU(n)\ltimes X$ has order exactly $n$ \cite[Thm.~1, Thm.~8]{paperA}
(\texttt{thmG\_general.py}). Weyl context families and their commutator-carry invariant
$Q$ are as in \cite{paperB}. Throughout, ``verified'' means: reproduced by the named
script in the companion repository \cite{repo}; expected one-line outputs are pinned in
\texttt{verification/INDEX.md}.

The repaired slogan is: \emph{the class detects the AvN sector of contextuality, no more
and no less}. The next three results make each clause exact.

\section{The equivalence on the AvN sector}
Fix $d\geq 2$, $m\geq 1$, and a Weyl context family $F$ whose contexts are commuting
triples multiplying to a central element. Write the value-assignment system of $F$ as
$A\gamma \equiv s \pmod d$, where $A$ is the context--observable incidence matrix and
$s$ collects the matrix-computed context phase exponents under the pinned Weyl
convention.

\begin{theorem}[Equivalence, AvN sector]\label{thm:equiv}
$F$ admits no noncontextual value assignment over $\mathbb{Z}_d$ if and only if some
cycle of $F$ carries odd anomaly self-pairing $Q$.
\end{theorem}

The criterion direction is \cite[Thm.~4]{paperB}. The solvability reformulation above is
verified two-directionally in \texttt{holonomy\_vs\_solvability.py}: over $1{,}500$ random
families at $d=2$ and $600$ at $d=4$, unsolvability of $A\gamma\equiv s$ (decided by
exhaustive dual search) coincided with the odd-$Q$ cycle test in every instance, with the
Peres--Mermin square ($d=2$) and the certified $d=4$ family of \cite{paperB} pinned as
positives so that both truth values of the equivalence are exercised at both $d$. The
achievable obstruction values themselves are classified as $\{0,d/2\}$ for even $d$ and
$\{0\}$ for odd $d$ \cite[Thm.~1]{paperB}.

\section{Both containments are strict}
\begin{proposition}[Nontrivial class, zero AvN]\label{prop:gap1}
On $\mathbb{Z}_3^2$ the Heisenberg cocycle $\omega(u,v)=\zeta_3^{\langle u,v\rangle}$ and
its square are coboundaries for none of the $3^9=19{,}683$ cochains; the class has order
exactly $3$. Simultaneously, every cycle of every sampled qutrit Weyl family has even
self-pairing, and the achievable AvN value at $d=3$ is $0$.
\end{proposition}
Verified exhaustively in \texttt{d3\_gap\_certificate.py}. Hence a nontrivial canonical
class does not imply operational (AvN) contextuality: the weakening adopted for the
gerbe-theoretic program \cite{cmp} is necessary, not cautious.

\begin{proposition}[Contextuality, zero class shadow]\label{prop:gap2}
The KCBS pentagon in $d=3$ attains quantum value $\sqrt{5}\approx 2.2361$ against the
exhaustive classical bound $2$, while the state-independent class shadow at $d=3$ is $0$.
\end{proposition}
Verified in \texttt{kcbs\_converse.py}. Hence state-dependent contextuality can be
invisible to the class: the converse direction of the naive slogan also fails.

\section{The Wigner--friend loop is level-split}\label{sec:wf}
Consider the friend-qubit loop $T \to T_x \to T_y \to T$ generated by
$U_1=\mathbb{1}\otimes H$, $U_2=\mathbb{1}\otimes S$.

\begin{observation}[Section dependence]\label{obs:wf}
The projective (Pancharatnam) holonomy of the loop is $+1$: the two eigenvector strands
carry Bargmann phases $e^{\pm i\pi/4}$ (the Bloch-octant $-\Omega/2$ law), whose product
is $1$, gauge-invariantly. The determinant-section holonomy is exactly $-i$ for every
special-unitary closing arrow; $-i=\tau^{-1}$ at $d=2$, a $\tau$-level invariant of the
2-adic tower, with section-independent square $-1$.
\end{observation}
Verified in \texttt{wf\_loop\_holonomy.py} (200 gauge randomizations; 500 closers). Any
physical or interpretational claim about ``the'' holonomy of this loop must state the
level; the two answers differ.

\paragraph{Hardware predictions.} The strand phases $\pm 45^\circ$ and loop product $+1$
are interferometrically measurable with two-qubit Hadamard tests; on the same device the
Peres--Mermin loop yields context signs $(+,+,+,+,+,-)$, i.e.\ ray-level holonomy $-1$.
One machine, two loops, opposite holonomies, both sharp.

\emph{Measured} (\texttt{ibm\_fez}, 2026-07-09, job \texttt{d986q62f47jc73a7hm2g}; registry
record \texttt{hardware/results/ibm\_fez\_holonomy\_20260709.json}): strand phases
$+45.88^\circ$ and $-48.81^\circ$ against targets $\pm45^\circ$ (per-strand statistical error
$\approx1.1^\circ$; the $\sim3.5\sigma$ offset on the $X$-prepared strand is a plausible
preparation systematic), loop phase $-2.93^\circ\pm1.55^\circ$ against target $0^\circ$. The
ray-level holonomy of the Wigner--friend loop is thus measured to be $+1$; the value $-i$
(loop $-90^\circ$) is excluded at this level by $\approx 56\sigma$, in accordance with
Observation~\ref{obs:wf}: $-i$ is a section-level, not interferometric, invariant. On the
same device the Peres--Mermin loop gave context products
$(+0.750,+0.835,+0.722,+0.780,+0.824,-0.700)$, sign pattern $6/6$ as predicted, and
$S=4.6125\pm0.0173$: a $35.4\sigma$ violation of the exact deterministic noncontextual bound
$4$ (the run-time script carried the loose $n_c-1=5$ bound, since replaced in
\texttt{qkernel} by the exhaustive tight bound), consistent with the archived
\texttt{ibm\_fez} value $4.558\pm0.018$.

\section{The residual problem}
\begin{openproblem}\label{op}
Characterize, in invariants of the context stack $[X/\PU(n)]$ refined beyond the
canonical $H^2$ class, the property ``some state has positive contextual fraction on
$F$''. Equivalently: identify the cohomological home of KCBS-type (state-dependent,
class-invisible) contextuality, and prove or refute that a sheaf-theoretic refinement
restores a two-directional equivalence.
\end{openproblem}
Propositions~\ref{prop:gap1}--\ref{prop:gap2} show no formulation at the level of the
canonical class alone can succeed; Theorem~\ref{thm:equiv} shows the AvN sector is
already optimally captured there.

\begin{lemma}[No fully covariant qubit-pair nets]\label{lem:nonet}
There is no two-qubit quasi-probability frame, covariant under all Pauli displacements,
whose every Lagrangian line-sum is a stabilizer projector: the $15$ Lagrangian sign
conditions have rank $10$ over $\mathbb{F}_2$ and are inconsistent, with minimal
certificate a Peres--Mermin-type magic configuration (nine observables, each in exactly
two of six contexts, $\varepsilon$-product $-1$). The order-$2$ class thus obstructs the
global frame itself (\texttt{ghw\_net\_necessity.py}).
\end{lemma}

Genuine Gibbons--Hoffman--Wootters nets evade Lemma~\ref{lem:nonet} by retreating to
$\mathrm{GF}(4)$ translation covariance; sweeping all $1024$ of them
(\texttt{gf4\_net\_necessity.py}) eliminates global frames in both remaining directions:
a state with contextual fraction $0.3462$ on the trivial-class family admits a
nonnegative representation in some net (necessity fails), while $T\otimes T$ is negative
in every net (best case $-1/4$) with contextual fraction $0$ (sufficiency fails); all
$60$ stabilizer states are nonnegative in some net (construction sanity). Together with
the $2^{15}$-frame sweep (\texttt{net\_robust\_negativity.py}), every global-frame
formulation of the even-$d$ classifier is refuted by explicit certificate; the
scenario-paired formulation below is not merely the surviving candidate but the forced one.

The constructive half now exists (\texttt{paired\_frame\_construction.py}). For a
closed-triple Weyl family $F$ at $d=2$, the quantum model $e(\rho)$ is a bijection of the
characteristic vector $c(\rho)=(\mathrm{tr}\,\rho T_v)_{v\in O(F)}$, and every
noncontextual model is supported on the solution set $L(F)$ of the family's sign system
(parity-violating outcomes have $e\equiv 0$). Hence:

\begin{theorem2}[Codeword-polytope classification; machine-validated]\label{thm:pair}
For closed-triple Weyl families at $d=2$,
$\;\mathrm{CF}(F,\rho)=0 \iff c(\rho)\in\mathrm{conv}\,L(F)$, the
$\varepsilon$-shifted codeword polytope of the family's $\mathbb{F}_2$ kernel; and
$L(F)=\varnothing$ exactly on the AvN sector, so one statement covers both sectors.
On the two trivial-class Peres--Mermin subfamilies the polytope has exactly $24$ facets,
all triangles $\pm c_u\pm c_v\pm c_w\le 1$; the classification holds $40/40$ on both,
while the canonical visible frame alone ($15/40$, $20/40$) and the induced CHSH squares
($31/40$, $29/40$) are each insufficient --- negative results retained as certificates.
\end{theorem2}

The pairing of Conjecture~\ref{conj} is thus, in the even-$d$ closed-triple case,
membership in a polytope cut out by the same $\mathbb{F}_2$ kernel that classifies the
AvN sector; what remains open is general $d$, mixed families, and the written proof of
the facet structure.

\subsection*{The ghost facet theorem and a closed formula}
The facet structure now has a proof and yields an exact formula
(\texttt{ghost\_facet\_theorem.py}). Let $F$ be Peres--Mermin minus one context $C^*$
with sign $\varepsilon^*$. \emph{Bijection lemma:} for a closed triple with
$T_uT_v=\varepsilon_C T_w$, the joint distribution is supported on
$s_us_vs_w=\varepsilon_C$ where $e_C(s)=\tfrac14(1+\sum_i s_i c_i)$; hence
$e(\rho)\leftrightarrow c(\rho)$ affinely. \emph{Support lemma:} noncontextual models
put no mass on parity-violating outcomes, so they are mixtures over $L(F)$, and a
mixture's model depends only on its mean, giving Theorem~\ref{thm:pair}.
\emph{Two-tetrahedra lemma:} every $\lambda\in L$ has ghost parity
$\prod_{v\in C^*}\lambda_v=-\varepsilon^*$ (the product of the five remaining signs,
$=-\varepsilon^*$ because the six Peres--Mermin signs multiply to $-1$ --- the class),
while every quantum state's ghost coordinates lie in the opposite tetrahedron of parity
$\varepsilon^*$ (since $T_1T_2T_3=\varepsilon^*\mathbb 1$ on $C^*$). Consequently, for
$\sigma$ with $\prod\sigma_i=\varepsilon^*$, every $\lambda\in L$ disagrees with
$\sigma$ in an odd number of positions, so $\sum_i\sigma_i\lambda_{v_i}\le 1$: these
four inequalities are valid, tight on $12$ of the $16$ vertices, and are the only
quantum-operative facets; the $20$ in-context facets are the probability positivity
conditions $e_C(s)\ge 0$, never violated by quantum states, and the opposite-parity
ghost directions are unconstrained ($L$ reaches $3$). The classification therefore
collapses to a formula:
\[ \mathrm{CF}(F,\rho)\;=\;\max\Bigl(0,\ \tfrac{S^*(\rho)-1}{2}\Bigr),\qquad
   S^*(\rho)=\max_{\prod\sigma=\varepsilon^*}\ \sum_{v\in C^*}\sigma_v\,
   \mathrm{tr}(\rho T_v), \]
verified to $10^{-16}$ across $60$ states on each of three deletions. The upper bound
$\mathrm{CF}\le(S^*-1)/2$ follows analytically from the validity above; equality now
follows from the \emph{Completion Lemma}: for each operative pattern $\sigma$ the point
$y_\sigma$ with ghost coordinates $\sigma$ and all others zero is a valid
parity-supported no-signalling model (every remaining context contains at most one ghost
observable, so its probabilities are $(1\pm1)/4$ or $1/4$), attaining $\sigma\cdot c=3$;
Abramsky--Barbosa--Mansfield duality then gives
$\mathrm{CF}=\max_\sigma(S_\sigma-1)/(3-1)$ exactly. The $d=2$ theorem is thus proved
end to end. Corollary: joint eigenstates of $C^*$ with an operative
pattern attain $S^*=3$ and hence $\mathrm{CF}=1$ --- for the deleted odd column these
are the Bell states, so \emph{maximal contextuality survives the deletion of the very
context that carried the class}: the class does not disappear, it flips the ghost
tetrahedron.

\subsection*{The ghost at $d=4$: the tower appears}
The mechanism lifts to the first nontrivial tower level
(\texttt{d4\_ghost\_tower.py}), on the certified $d=4$ magic square (six contexts, nine
observables of mixed orders $2$ and $4$, each observable in exactly two contexts, phase
exponents summing to $d/2$ --- the class value). Deleting any context $C^*$ leaves a
solvable system with $|L|=64$ spectral assignments, and the ghost value
$\sum_{v\in C^*}\lambda_v$ is constant on $L$ with
\[ \gamma - s^{*} \;\equiv\; d/2 \pmod d : \]
\emph{the classical ghost torus is rotated from the quantum one by exactly the top
2-adic layer} $\omega^{d/2}=-1$, for both a pure order-$2$ ghost and a mixed
order-$2/4$ ghost --- and the $d=2$ opposite-tetrahedra lemma is the same law, since
there too the offset was $d/2$. Every joint eigenstate class of the deleted context
attains $\mathrm{CF}=1.0000$ exactly ($4$ and $8$ classes respectively), and
Haar-random states populate the sector (sampled $\mathrm{CF}$ up to $0.21$).
Two negative results are pinned honestly: noncontextual assignments must be restricted
to spectral values (non-spectral assignments hit zero-probability outcomes; the
unrestricted relaxation under-computes $\mathrm{CF}$), and the $d=4$ law is \emph{not}
a function of the Hermitian order-$2$ shadows of the ghost alone --- the Hermitianized
$d=2$ formula predicts zero precisely where random states reach $\mathrm{CF}=0.138$.
The exact $d=4$ classification is polytope membership in the full moment space of
Theorem~\ref{thm:pair}. Two further results pin its structure
(\texttt{d4\_moment\_facets.py}).

\emph{Ghost-law lemma (now proved).} Since every observable lies in exactly two
contexts, summing the remaining constraints over any $\lambda\in L$ gives
$\gamma \;=\; \sum_{\mathrm{rem}} s \; - \; 2T \pmod d$, with
$T=\sum_{v\notin C^*}\lambda_v \bmod 2$; using $\sum_{\mathrm{all}} s = d/2$ (the
class) this reads $\gamma-s^{*} = d/2 - 2(s^{*}+T)$. $T$ is constant on $L$ and
$s^{*}+T$ is even for every one of the six deletions (machine-verified), so
$\gamma-s^{*}=d/2$ always: the tower ghost law is an identity plus a verified parity
sub-lemma.

\emph{$\tau$-necessity (twelve explicit certificates).} For each sampled state with
$\mathrm{CF}>0$, the \emph{optimal} separating functional restricted to the
even-character subspace --- all order-$2$ shadow data, i.e.\ everything visible below
the top 2-adic layer --- achieves gap $\le 0$: the even projection of the quantum moment
vector lies \emph{inside} the even projection of the classical polytope. Hence the
$d=4$ state sector is invisible at the order-$2$ layer, and the $\tau$-level moments are
necessary for every operative facet. The 2-adic tower is not bookkeeping; it is
load-bearing for the state sector, exactly as the even-$d$ half of
Conjecture~\ref{conj} asserts. The odd-sector facet structure is now partially decoded
(\texttt{d4\_odd\_sector\_facets.py}). The polytope has affine dimension $33$; sparse
certificate mining shows the operative geometry is two-tiered. \emph{Tier one} consists
of triangle inequalities on \emph{closed $\tau$-twisted order-$2$ triples}: every
$3$-sparse certificate coordinate ($21/21$) is an order-$2$ operator reached through
odd character indices of its host context --- the shadow wearing the top-layer cocycle
phase --- and the recurring triples are the deleted context $C^*$ itself (the
$\tau$-twisted ghost) and a virtual closed triple
$\{(0,2,2,0),(2,0,0,2),(2,2,2,2)\}$. This tier classifies $32/40$ sampled states.
\emph{Tier two}, required by the remaining contextual states, involves genuinely
order-$4$ operator moments --- the second level of the tower --- and is now pinned
exactly (\texttt{d4\_tier2\_catalogue.py}): the extracted supporting inequalities have
tight-vertex rank exactly $32$ (true facets of the $33$-dimensional hull), their normals
snap with \emph{zero} error onto the $\tfrac{1}{2\sqrt2}\mathbb Z$ grid with a single
coefficient magnitude $3/(2\sqrt2)$, and their bounds are the integers $6$ and $7$ ---
equivalently, in unit-coefficient normalization, $4\sqrt2$ and $14\sqrt2/3$:
\emph{Tsirelson's $\sqrt2$ sits in the tier-two bounds}, with $\tau=e^{i\pi/4}$-quantized
phases in the normals. Together, tier one and the exact tier-two facets classify $40/40$ states of the
construction sample; a subsequent fresh-sample audit shows this is \emph{not yet the
complete family} --- new weak-band contextual states require further tier-two orbit
members. Closing the extracted seeds under the two proven vertex symmetries (kernel
translations and global conjugation, $\sim\!1152$ facets) recovers most but not all
fresh states, and per-context outcome relabelings are provably \emph{not} symmetries
(each observable lives in two contexts). The deleted family's automorphism group is now derived: the family is $K_{3,3}$
(contexts as nodes, observables as edges), and the automorphisms fixing the deleted
context form the twelve graph automorphisms dressed by a $64$-torsor of spectral
shifts, times conjugation --- order $1536$, each element provably a bijection of the
$64$ vertices, acting on moment space by exact monomial rotations. The census is now \emph{complete and exact}. The polytope --- $64$ integer
$\{-1,0,1\}$ points in faithful $33$-dimensional integer coordinates --- has exactly
$\mathbf{23{,}256}$ facets in $\mathbf{61}$ symmetry classes, derived twice by
independent methods that agree class-for-class: an adjacency decomposition (breadth-first
closure of the facet--ridge graph modulo the derived group, every pivot certified by
exact integer nullspace recovery through rank-$32$ tight sets and exact validity on all
vertices), and a direct exact enumeration of the full polytope by \textsc{Normaliz}.
Tight-vertex counts run $36$--$60$; orbit sizes run $4$ to $768$, so the census contains
both highly symmetric facets (stabilizers of order $192$) and free ones. Since quantum
moment vectors lie in the affine hull (residual $\sim 10^{-14}$), the catalogue is the
classification theorem in operational form: $\mathrm{CF}(\rho)>0$ if and only if
$\mu(\rho)$ violates one of the $23{,}256$ inequalities; on fresh sampled states the
facet classifier agrees with the LP $40/40$, as it must. Two method notes are pinned in
the suite: \texttt{lrs} wedges reproducibly on several of these tight-set polytopes
(stalling at identical partial output under different budgets) while \texttt{cdd} and
\textsc{Normaliz} succeed; and an intermediate in-walk count of ``$121$ classes''
was a bookkeeping duplication (two key encodings double-storing each orbit) exposed
only by the independent cross-check --- the cross-check is not optional. The tower
appears three times over: as the ghost rotation $d/2$, as the stratification of
operative facets by operator order, and as the $\mathbb Z[1/\sqrt2]$ arithmetic of the
tier-two facet data. Open: the no-signalling-normalized closed formula, mixed families,
and the arithmetic profile of all $61$ classes.

\begin{conjecture}[State-sector pairing]\label{conj}
For odd prime $d$, the state-dependent sector is classified by the pairing of the state
against the order-$d$ class through the Weyl transform: a state yields positive
contextual fraction on stabilizer measurement scenarios if and only if its parity-frame
quasi-distribution is negative --- the Howard--Wallman--Veitch--Emerson
contextuality--negativity correspondence \cite{hwve}, recast as a stack pairing. For even
$d$, where no Gross frame exists, the corresponding role is played by the top layer of the
2-adic tower of \cite{paperB}.
\end{conjecture}

Computed anchors (\texttt{state\_sector\_probe.py}): $\mathrm{CF}(\text{PM})=1.0000$ for
every state (AvN); $\mathrm{CF}(\text{KCBS})=0.4721=2\sqrt5-4$ at the optimal state with
class shadow $0$; at $d=3$ the parity-frame quasi-distribution is nonnegative on all
twelve stabilizer pure states (exhaustive) and attains $-1/3$ on the strange state.
Notably, the $\tau$-convention frame of \cite{paperB} exhibits $-1/18$
pseudo-negativity on stabilizer states at $d=3$: $\tau$ is the even-$d$ object, and the
odd/even divide of \cite{paperA} resurfaces at the level of quasi-probability frames.

\appendix
\section{Claims-to-verification table}
\begin{center}\small
\begin{tabular}{lll}
\toprule
Claim & Script & Expected one-liner\\
\midrule
Thm.~\ref{thm:equiv} & \texttt{holonomy\_vs\_solvability.py} & $0$ mismatches; pinned PM/cert4 agree\\
Prop.~\ref{prop:gap1} & \texttt{d3\_gap\_certificate.py} & order $3$; AvN $\{0\}$; gap certified\\
Prop.~\ref{prop:gap2} & \texttt{kcbs\_converse.py} & $2.23607>2$; converse gap certified\\
Obs.~\ref{obs:wf} & \texttt{wf\_loop\_holonomy.py} & strands $\pm\pi/4$; $\det$-holonomy $-i$\\
PM anchor & \texttt{pm\_gauge\_invariance.py} & product $-1$ over all $512$ gauges\\
\bottomrule
\end{tabular}
\end{center}

\begin{thebibliography}{9}
\bibitem{paperA} M.~Flores Gordillo, \emph{The Canonical Class of the MASA Context Stack
and the Peres--Mermin Obstruction}, working draft, 2026.
\bibitem{paperB} M.~Flores Gordillo, \emph{The Obstruction Spectrum of Weyl Context
Families and the 2-adic Tower}, working draft, 2026.
\bibitem{cmp} M.~Flores Gordillo, \emph{Contextual Holonomy, Lifting Bundle Gerbes, and
the Geometric Structure of Quantum Contextuality}, in preparation.
\bibitem{hwve} M.~Howard, J.~Wallman, V.~Veitch, J.~Emerson, \emph{Contextuality supplies the magic for quantum computation}, Nature 510, 351 (2014).
\bibitem{repo} \texttt{github.com/manuflog/contextuality-obstructions}, verification
suite, 2026.
\end{thebibliography}
\end{document}
